\chapter{Ноды, независимость, остановка}

\section{Две формы клиента}

Рабочая нода это агентский цикл контрибьютора. Протокол допускает две формы клиента, и выбор между ними определяется тем, кому доверяет проект.

\term{Публичный рой} использует минимальный runtime, план которого содержит ровно две привязки: удалённый shell среды и подпись ревизий. Runtime не поставляет инструментов сам по себе, поэтому донорский профиль есть свойство конструкции, а не настройка: расширить его нельзя, не изменив digest плана, а digest плана сверяет среда на каждом вызове.

\term{Доверенная организация}, компания или кооператив, может вместо этого подключить среду как сервер инструментов внутри родного харнеса вендора. Донорский профиль там держится политикой организации: у агента остаются локальные инструменты, и проект доверяет тому, что организация их контролирует.

\begin{whybox}[Почему для публичного роя не подходит харнес вендора]
Харнес вендора всегда несёт локальный shell и локальные файлы. Задача от незнакомого проекта есть недоверенный ввод, и инжект в ней получает в распоряжение машину контрибьютора. Ограничить харнес снаружи нельзя надёжно; можно только не поставлять лишнего. Минимальный runtime с двумя привязками решает это по построению. Цена: подписки на модели можно использовать только там, где условия вендора допускают сторонний runtime; иначе источником ёмкости становятся ключи API и локальные модели.
\end{whybox}

\term{Решение 28 сентября 2026 года.} Собственный агентский цикл отложен: и оркестраторы, и рабочие ноды работают на существующих харнесах (Claude Code, Codex, Muse, харнес над локальной моделью), и независимость от харнеса есть принцип дизайна (curator-keeper, \texttt{spec/keeper.md} v0.5.1, §7.7). Минимальный runtime откладывается вместе с циклом. Пока его нет, донорский профиль публичной формы держится переопределением инструментов существующего харнеса; именно этот путь спецификация главного цикла называет возможным первым путём к ограниченным рабочим агентам. Это слабее свойства конструкции, потому что харнес по-прежнему несёт локальный shell и файлы, а профиль держится ровно настолько, насколько надёжно переопределение. Врезка выше остаётся верной как довод о том, чего переопределение должно достичь.

Обе формы говорят с проектом одним протоколом ноды: присоединение по ключу, heartbeat, назначение, сессии удалённого исполнения, представление ревизии, ревью, выход.

\section{Ёмкость любого происхождения}

Ресурс роя есть inference-ёмкость, а не подписка. Источники: подписки, где условия это допускают; ключи API; локальные модели на собственном железе контрибьютора. Разнообразие происхождения не только расширяет предложение, но и снижает корреляцию ошибок ревью: два ревьювера на разных семействах моделей пропускают разное.

\begin{whybox}[Что на самом деле отдаёт контрибьютор]
Folding@home работает потому, что простаивающий процессор ничего не стоит владельцу. Квота подписки не простаивает, её хотят потратить сами. Простаивает только то, что сгорит: неиспользованный остаток окна квоты перед сбросом. Это и есть главный ресурс публичного роя. Heartbeat сообщает остаток, окно сброса и класс стоимости; планировщик берёт истекающую ёмкость первой. Нода ограничивает расход на проект за окно, чтобы рой не съел то, что владелец хотел оставить себе.
\end{whybox}

\section{Подпись на машине владельца}

Ключ ноды никогда не оказывается на среде. Среда возвращает хеши состояний и вывода; провайдер на машине ноды подписывает ревизию и receipts (инвариант I6). В runtime это устроено так: ядро, которое обрабатывает вывод модели, не имеет доступа к ключам вовсе; ими владеет брокер, и подпись есть операция брокера, а не инструмента модели. Модель просит подписать хеш, брокер подписывает то, что видит в журнале задачи, и не подписывает ничего другого. Собственный runtime отложен (раздел 8.1), поэтому на существующих харнесах роль брокера должен играть компонент на машине владельца вне процесса модели. Ближе всего к этому хост сессий: по waggle он уже держит ключ сессии и подписывает за агента, который ключа не видит (waggle, \texttt{spec/waggle.md} v5.2, §4.2).

\begin{center}
\begin{tikzpicture}[node distance=8mm]
\node[boxa,text width=34mm] (core) {Ядро агента\\[1pt]\footnotesize обрабатывает вывод модели, ключей нет};
\node[boxg,text width=34mm,right=17mm of core] (broker) {Брокер\\[1pt]\footnotesize держит ключи, подписывает хеши};
\node[boxw,text width=34mm,right=17mm of broker] (env) {Среда\\[1pt]\footnotesize возвращает хеши, ключей не видит};
\draw[arr] (core) -- node[lbl,above] {подпиши\\хеш $R$} (broker);
\draw[darr] (broker) -- node[lbl,above] {команды,\\хеши} (env);
\end{tikzpicture}
\end{center}

\section{Донорский профиль и инжект}

Задача от проекта есть недоверенный ввод для агента. В публичной форме клиента план из двух привязок означает, что инжект в задаче не имеет локальных инструментов, к которым мог бы обратиться: он может только исполнять команды в санбоксе среды, куда и так был направлен, и просить подписать хеши, которые брокер сверяет с журналом. Бюджеты на проект за окно ограничивают ущерб от задачи, которая просто жжёт ёмкость.

\section{Конфиденциальность}

Рабочая нода видит содержимое своего санбокса и отправляет его провайдеру своей модели. Протокол не скрывает код от контрибьюторов и их провайдеров, и это нужно понимать до присоединения. Манифест проекта объявляет допустимых провайдеров и условия хранения; проект, которому это неприемлемо, ограничивает контрибьюторов или провайдеров, и это его политика, а не свойство протокола.

\section{Три оси независимости}

Два владельца на одной модели с одним харнесом ошибаются одинаково. Поэтому независимость определяется по трём осям:

\begin{itemize}
\item административный домен, против сговора и множественных личностей;
\item семейство модели, против общих слепых зон;
\item тулчейн, против одинаковых ошибок среды и анализаторов.
\end{itemize}

Независимость есть политика проекта с проверяемыми основаниями, а не пожелание. Проект объявляет, какие оси обязательны, и для каждого сведения о ноде говорит, к какому из трёх видов оно относится: подтверждается (ключ ноды и способ оплаты ёмкости подтверждают владельца в меру доверия к ним), принимается на доверии (объявленное семейство модели) или служит эвристикой. Различие строк в объявлении не доказывает различия владельцев.

\begin{warnbox}[Политика не ослабляется сама]
Если подходящих участников нет, задача ждёт, эскалируется либо получает отдельно авторизованное изменение условий, записанное в журнал. Реализация, которая молча снижает требования, когда некому ревьюить, ломает инвариант I5 в самом уязвимом месте.
\end{warnbox}

Отдельно различаются независимость ревью и независимость исполнения проверок. Третья нода, запустившая проверки в той же среде, даёт разделение ролей, но не независимое подтверждение честности среды; его даёт только верификатор на другой машине. Формула вероятности пропуска в главе 9 предполагает независимость ошибок; три оси делают ошибки менее коррелированными, но не независимыми, и это надо помнить при чтении любых чисел.

\section{Идентичность и привязка ключей}

Ось владельца из объявленной становится проверяемой через доказательства идентичности. Это не отдельная подсистема, а способ заполнить графу подтверждается в политике независимости.

\subsection{Двустороннее доказательство}

Владелец ноды имеет долгоживущий \term{ключ идентичности} и доказывает его принадлежность внешней сущности с двух сторон: заявлением я есть X у провайдера P, подписанным ключом, и артефактом на стороне провайдера, в котором лежит отпечаток ключа. Провайдеры суть \term{резолверы} с одним интерфейсом: по строке идентичности вернуть набор отпечатков ключей и хеш свидетельства.

\begin{likbez}[{Прецеденты: Keybase, DNS, well-known}]
Схема с заявлением у себя и отпечатком у провайдера известна по Keybase, где пользователь доказывал владение аккаунтами соцсетей, доменом и репозиторием. Публикация ключей в DNS стандартизована в DANE и в записях OPENPGPKEY; Bluesky привязывает handle к идентичности TXT-записью домена; Nostr публикует ключи по well-known адресу. Ничего из этого не нужно изобретать, нужно выбрать набор резолверов и определить, как результат их работы попадает в журнал.
\end{likbez}

Стартовый набор резолверов: TXT-запись домена под DNSSEC для организаций; well-known по HTTPS; ключи аккаунта на хостинге кода, потому что контрибьюторы это разработчики и ключ у них уже опубликован. Публикации в социальных платформах как необязательный класс с меньшим весом. Системы имён с блокчейном как плагины позже, если понадобятся. Протоколу важен только результат резолвера и его свидетельство.

\subsection{Подтверждение как запись журнала}

Резолв есть внешний вызов и недетерминирован между валидаторами: DNS не разошёлся, лимит API, сеть. Поэтому идентичность считается подтверждённой не потому, что кто-то проверил, а потому, что $t$ валидаторов независимо получили и записали одинаковый хеш свидетельства в одном окне. Расхождение даёт не подтверждено, без исключений. Это то же разделение подписанного свидетельства и независимо установленного факта, что и в главе 9, и оно ложится в секвенсер с асинхронными подписями без нового механизма.

Свидетельство хранится в журнале с датой. Пост удалили, API закрыли, домен временно недоступен: зафиксированное наблюдение остаётся, а политика проекта задаёт каданс переподтверждения. Первое увиденное свидетельство пинится: смена ключа в доказательстве есть событие с записью, а не молчаливое обновление.

\subsection{Ключ идентичности и ключи нод}

Ключ идентичности долгоживущий, он лежит в доказательствах и меняется редко. Ключи нод и устройств короткоживущие и сертифицируются ключом идентичности записью в журнале с началом и концом действия. Нода ротирует ключи, не переделывая доказательств; репутация привязана к идентичности и переживает ключи и проекты. Ротация через перекрытие и отзыв удалением у нас уже есть в реестре Curator для кор-ключей, здесь тот же набор пинов уровнем выше. Отзыв прост: убрать отпечаток из TXT-записи или well-known, и первый валидатор, не нашедший его, записывает расхождение.

\subsection{Подпись коммитов и сообщений}

Ревизии подписываются ключом ноды в формате SSH-подписи git. Файл допустимых подписантов для среды выводится из журнала: ключ допустим, если сертифицирован подтверждённой идентичностью и назначен на задачу. Хостинг кода покажет подпись как проверенную побочно, но проверяющим является журнал, а не платформа. Сообщения протокола уже подписаны ключами нод; добавляется только цепочка до ключа идентичности.

\subsection{Что это покупает и чего нет}

Покупает: проверяемую независимость владельцев (два домена или две сертифицированные идентичности, а не две строки в объявлении); рост цены множественных личностей, потому что домен или старый аккаунт дороже ключа за секунду; переносимую репутацию; отзыв; подотчётность людям, а не отпечаткам. Симметрично для проектов: ключи кор-команды публикуются под доменом организации, и контрибьютор при присоединении проверяет подписантов манифеста через домен.

\begin{warnbox}[Чего доказательство не удостоверяет]
Контроль над доменом не есть честность, аккаунт можно купить, DNS без DNSSEC угоняется. Доказательство удостоверяет контроль, а не поведение. Поэтому политика назначает вес по классу: домен под DNSSEC выше социального поста, а псевдонимная нода допускается с меньшим весом и большей глубиной ревью, но не запрещена. Запрет псевдонимов отрезал бы часть предложения ёмкости, не купив безопасности.
\end{warnbox}

В план MVP слой идентичности не входит: журнал уже умеет хранить ключи и доказательства как записи, и резолверы добавляются, не трогая инвариантов. Это первый релиз после MVP.

\section{Репутация}

Нода идентифицируется ключом и объявляет три оси. Стоимость создания ноды равна стоимости ёмкости, что даёт естественный барьер против множественных личностей, но нода на дешёвой модели тоже нода, поэтому барьер дополняется репутацией.

После каждой задачи проект выдаёт ноде подписанную аттестацию: принято, отклонено, вердикт ревью совпал с итогом или нет. Репутация определяет глубину проверки, как адаптивная репликация в BOINC: новой ноде назначаются два или три независимых ревьювера, ноде с историей один. Она определяет и вес при выборе в ревьюверы. Пропорциональное вознаграждение при полной устойчивости к множественным личностям недостижимо, поэтому репутация снижает риск, а не устраняет его.

Аттестации переносимы между проектами и проверяются офлайн. Проект может выдавать их собственным подставным нодам, поэтому чужие аттестации взвешиваются репутацией выдавшего проекта.

\section{Остановка}

Любая кор-нода выпускает сертификат экстренной остановки: подпись одного ключа, эпоха и монотонный номер против повтора. Среда по нему немедленно переходит в безопасный режим: не открывает санбоксы, не проводит операции раннера, не продвигает основную ветку. Позднее остановка упорядочивается журналом. Возобновление требует сертификата кворума уровня.

\begin{whybox}[{Почему одна подпись останавливает, а кворум возобновляет}]
Асимметрия намеренная: остановка дешёвая и безопасная, продолжение дорогое. В системе, где порча основной ветки хуже простоя, целостность важнее доступности. Нода, злоупотребляющая остановкой, исключается записью о составе, которая сама проходит кворум. Скомпрометированный оркестратор не может заблокировать остановку, потому что она не проходит через него, а откат журнала на версию до остановки невозможен по инварианту I8.
\end{whybox}
